\section{Conclusion}
\label{sec:rewrite-conclusion}

The endogenous-growth model shows when improvements in AI efficiency generate
a temporary increase in growth and when they support a permanently
faster-growing economy. In the labor-bottleneck regime, the disappearance of
the growth premium does not imply that earlier gains in output and wages
disappear. In the AI-dominated regime, real wages grow permanently faster,
but output per worker grows faster still. Workers become better paid while
receiving a declining share of aggregate income.

These conclusions are conditional on the production and research technologies.
The link between permanently faster wage growth and a vanishing labor share
applies to the long-run regimes I characterize with a finite upper bound on
AI efficiency. Without this bound, the unit-elastic benchmark can instead
sustain faster growth with a positive labor share. The benchmark also assumes
autonomous AI research. Introducing human research tasks and allowing the
attainable level of efficiency to evolve would test how these assumptions
shape the results. The research networks and bottlenecks in
\citet{davidsonetal2026} offer a starting point, but their conditions for
explosive growth do not establish equilibrium existence in this model.

\ifshowlegacysimulations
The transition can be as important as the long-run outcome. In the calibrated
AI-dominated path, labor completes half of its income-share decline after 50
model years while output and wages grow rapidly. This timing is illustrative:
research productivity is not estimated, and changing it produces faster or
slower equilibrium adjustments rather than merely relabeling time. The initial
technological distance matters as well. The simulations nevertheless reveal
signals before most of the decline in labor's income share occurs: effective AI
production services can grow rapidly, the interest rate can rise above the
labor-bottleneck benchmark, and a positive output--wage growth gap shows that
labor's income share is already falling. Starting farther from the terminal
regime can require the developer's owner to finance large temporary losses
against future monopoly rents.
\else
The numerical exercises are illustrative rather than joint empirical fits to
AI revenues, research expenditure, and price changes. Research productivity
is particularly uncertain, so I place more weight on the conditional
long-run results than on the dates of takeoff. The measurement difficulties
documented by \citet{korinekmckelvey2026}, including limited information on
the allocation of compute, also constrain calibration. Moreover, the
experiments assume an unanticipated activation of RSI and omit adjustment
costs. An explicit compute sector with investment frictions, as in the GATE
model of \citet{erdiletal2025gate}, would allow me to examine whether the
large initial reallocations survive constraints on expanding compute capacity.
\fi

Permanent monopoly also restricts the interpretation of prices and profits.
Together with the production technology, it can imply a substantial AI revenue
share under complementarity even when the AI production weight is small.
I therefore retain this weight as an illustrative parameter rather than an
estimate of the industry's current size. \citet{korinekvipra2025} discuss
both forces toward concentration and competition among AI developers;
their analysis does not imply permanent monopoly. Allowing entry and
competition would make it possible to study how lower markups and the
prospect of losing market leadership affect private RSI investment.
AI revenue would still need to be distinguished from net profit after
inference and research costs.

The aggregate labor input also excludes differences across occupations and
the creation of new tasks in which workers retain a comparative advantage.
A task-based extension could incorporate the displacement and reinstatement
mechanisms of \citet{acemoglurestrepo2019}, allowing labor's role to change
through the creation and automation of tasks as well as through changes in
relative input use.
The present model cannot determine which workers gain or lose during
the transition.

Finally, a declining labor share does not by itself determine household
inequality or workers' economic power. The representative household owns
capital and the AI developer, so I do not distinguish workers who rely on
wages from households that receive most of the profits. As
\citet{jones2026future} emphasizes, asset ownership and redistribution matter
for who shares in AI-driven prosperity. Heterogeneous ownership and a fiscal
sector would allow me to study how capital income and transfers modify
the distribution of gains described here through wages and factor shares.

\enlargethispage{\baselineskip}
The central question is whether human labor remains sufficiently essential for
wages to keep pace with the economy as a whole.
